\chapter{Results and Discussion}
\label{chap:results}

\section{Simulation Results}

The simulation results for the two controllers, Linear Quadratic Regulator (LQR) and Data-Driven Min-Max Model Predictive Control (DD-MPC), are presented and compared.

Figures~\ref{fig:lqr_sim_lean} and~\ref{fig:ddmpc_sim_lean} show the evolution of the lean angle ($\phi$) over time for the LQR and the Data-Driven MPC controllers, respectively.

%\begin{figure}[H]
%	\centering
%	\includegraphics[width=0.8\textwidth]{lqr_sim_lean.png}
%	\caption{Lean angle over time -- LQR Controller (Simulation)}
%	\label{fig:lqr_sim_lean}
%\end{figure}
%
%\begin{figure}[H]
%	\centering
%	\includegraphics[width=0.8\textwidth]{ddmpc_sim_lean.png}
%	\caption{Lean angle over time -- Data-Driven MPC Controller (Simulation)}
%	\label{fig:ddmpc_sim_lean}
%\end{figure}

\subsection{Comparison of Controllers}

Key performance metrics from the simulations are summarized in Table~\ref{tab:sim_performance}.

\begin{table}[H]
	\centering
	\caption{Simulation Performance Comparison}
	\label{tab:sim_performance}
	\begin{tabular}{|l|c|c|}
		\hline
		\textbf{Metric} & \textbf{LQR} & \textbf{Data-Driven MPC} \\ \hline
		Maximum lean angle deviation ($^\circ$) & 6.2 & 4.3 \\ \hline
		Stabilization time (s) & 2.8 & 2.2 \\ \hline
		Overshoot ($^\circ$) & 1.5 & 0.8 \\ \hline
	\end{tabular}
\end{table}

The Data-Driven MPC controller achieved:
\begin{itemize}
	\item Approximately $30\%$ reduction in maximum lean angle deviation.
	\item Around $20\%$ faster stabilization time.
	\item Lower overshoot compared to the LQR controller.
\end{itemize}

These results validate the theoretical expectation that Data-Driven MPC provides enhanced robustness against model uncertainties and disturbances.

\section{Experimental Results}

Experimental trials were conducted for both controllers on the real mini-bike platform.

Figures~\ref{fig:lqr_exp_lean} and~\ref{fig:ddmpc_exp_lean} present the experimental lean angle responses.

%\begin{figure}[H]
%	\centering
%	\includegraphics[width=0.8\textwidth]{lqr_exp_lean.png}
%	\caption{Lean angle over time -- LQR Controller (Experiment)}
%	\label{fig:lqr_exp_lean}
%\end{figure}
%
%\begin{figure}[H]
%	\centering
%	\includegraphics[width=0.8\textwidth]{ddmpc_exp_lean.png}
%	\caption{Lean angle over time -- Data-Driven MPC Controller (Experiment)}
%	\label{fig:ddmpc_exp_lean}
%\end{figure}

\subsection{Experimental Observations}

\begin{itemize}
	\item The Data-Driven MPC maintained stability even under larger initial lean disturbances.
	\item The steering control inputs exhibited smoother behavior compared to the LQR.
	\item Some minor deviations from the simulation results were observed, primarily due to sensor noise, actuation delays, and friction.
\end{itemize}

Table~\ref{tab:exp_performance} summarizes key experimental performance metrics.

\begin{table}[H]
	\centering
	\caption{Experimental Performance Comparison}
	\label{tab:exp_performance}
	\begin{tabular}{|l|c|c|}
		\hline
		\textbf{Metric} & \textbf{LQR} & \textbf{Data-Driven MPC} \\ \hline
		Maximum lean angle deviation ($^\circ$) & 7.1 & 5.0 \\ \hline
		Stabilization time (s) & 3.0 & 2.4 \\ \hline
		Overshoot ($^\circ$) & 2.1 & 1.2 \\ \hline
	\end{tabular}
\end{table}

\section{Discussion}

The results demonstrate that the Data-Driven Min-Max MPC approach provides superior stabilization performance compared to traditional model-based LQR control, both in simulations and real-world experiments.

Major findings:
\begin{itemize}
	\item \textbf{Improved Robustness:} Data-driven methods successfully handled system uncertainties without requiring explicit modeling.
	\item \textbf{Reduced Stabilization Time:} Faster recovery from initial disturbances was observed.
	\item \textbf{Lower Overshoot:} Smoother convergence toward upright position minimized oscillations.
\end{itemize}

\subsection{Challenges and Limitations}

Despite the promising results, some limitations were encountered:
\begin{itemize}
	\item \textbf{Sensor Noise:} MPU6050 noise impacted lean angle estimation at low speeds.
	\item \textbf{Control Saturation:} Steering servo saturation limited rapid response capability.
	\item \textbf{Offline Computation:} Due to STM32 limitations, online real-time optimization of the SDP was not feasible; a precomputed gain was used instead.
\end{itemize}

\subsection{Future Work}

Potential improvements and future research directions include:
\begin{itemize}
	\item Online adaptive Data-Driven MPC to incorporate real-time data updates.
	\item Implementation of rear-wheel steering stabilization for stationary balancing.
	\item Integration of higher precision IMUs to reduce noise.
	\item Hardware upgrades to allow full real-time online optimization on embedded platforms.
\end{itemize}

